%---------------------------------------------------------------------
% Part one: Surah Al-Baqarah. Seven questions.
%
% Every verse quoted here was actually set in one of the four papers.
% Kept deliberately concise -- 3 to 5 lines of tafsir per point, not
% a full essay -- to hold the page budget across fourteen questions.
%---------------------------------------------------------------------

\section*{حصہ اول --- سورۃ البقرہ}
\markboth{سورۃ البقرہ}{}

\begin{nukta}[سورت کا خاکہ]
\textbf{مدنی} سورت، قرآن کی \textbf{دوسری اور سب سے بڑی} سورت، \textbf{دو سو
چھیاسی} آیات، تیس رکوع، پہلے تین پارے۔ نام سورت کے اندر مذکور ''گائے''
(\textenglish{69--71}) کے واقعے سے ہے۔ مرکزی مضامین: عقائد کی بنیاد، احکامِ
شریعت (نماز، زکوٰۃ، روزہ، طلاق، وراثت)، بنی اسرائیل کا احتساب، اور منافقین کی
نشان دہی۔
\end{nukta}

%---------------------------------------------------------------------
\begin{sawal}{سوال \textenglish{1} \ \textnormal{\small(بہار \textenglish{2013}، سوال \textenglish{1})}}
سورۃ البقرہ کی روشنی میں خلافتِ آدم علیہ السلام کے تصور پر نوٹ لکھیے۔
\end{sawal}

\begin{hal}
\ayat{وَإِذْ قَالَ رَبُّكَ لِلْمَلَائِكَةِ إِنِّي جَاعِلٌ فِي الْأَرْضِ
خَلِيفَةً قَالُوا أَتَجْعَلُ فِيهَا مَن يُفْسِدُ فِيهَا وَيَسْفِكُ الدِّمَاءَ
وَنَحْنُ نُسَبِّحُ بِحَمْدِكَ وَنُقَدِّسُ لَكَ قَالَ إِنِّي أَعْلَمُ مَا لَا
تَعْلَمُونَ}{البقرہ، \textenglish{30}}
\tarjuma{اور جب تیرے رب نے فرشتوں سے فرمایا کہ میں زمین میں ایک خلیفہ بنانے
والا ہوں۔ انہوں نے عرض کیا: کیا تو زمین میں ایسے کو بنائے گا جو اس میں فساد
کرے گا اور خون بہائے گا، حالانکہ ہم تیری حمد کے ساتھ تسبیح کرتے اور تیری
پاکی بیان کرتے ہیں؟ فرمایا: میں جانتا ہوں جو تم نہیں جانتے۔}

\smallskip
\textbf{تشریح:} ''خلیفہ'' یعنی زمین میں اللہ کے احکام نافذ کرنے والا نائب ---
یہ اعزاز فرشتوں کو نہیں، انسان کو دیا گیا۔ فرشتوں کا سوال بدگمانی نہیں،
استفسار تھا؛ جواب میں اللہ نے اپنے علمِ محیط کی طرف اشارہ کیا۔ آگے آیت
\textenglish{31} میں آدم\,علیہ السلام کو \emph{اسماء} (علم) کی تعلیم دی گئی،
جو خلافت کی اصل بنیاد ثابت ہوئی --- علم ہی وہ چیز ہے جو فرشتوں پر انسان کی
فضیلت کا سبب بنی۔

\medskip
\textbf{خلاصہ۔} انسان زمین پر امانت دار بنایا گیا، مالک نہیں۔ خلافت کا حق
علم اور تابعداری سے مشروط ہے، محض تخلیق سے نہیں۔
\end{hal}

%---------------------------------------------------------------------
\begin{sawal}{سوال \textenglish{2} \ \textnormal{\small(بہار \textenglish{2022}، سوال \textenglish{1} اور سوال \textenglish{5})}}
سورۃ البقرہ کی روشنی میں منافقین کی صفات بیان کریں، اور منافق و کافر میں فرق
واضح کریں۔
\end{sawal}

\begin{hal}
سورت کا آغاز تین گروہوں سے ہوتا ہے: مومن (\textenglish{2--5})، کافر
(\textenglish{6--7})، اور منافق (\textenglish{8--20}) --- اور منافقین کا بیان
سب سے \textbf{طویل} ہے، کیونکہ ان کا خطرہ چھپا ہوا ہے۔

\ayat{وَإِذَا قِيلَ لَهُمْ لَا تُفْسِدُوا فِي الْأَرْضِ قَالُوا إِنَّمَا
نَحْنُ مُصْلِحُونَ. أَلَا إِنَّهُمْ هُمُ الْمُفْسِدُونَ وَلَٰكِن لَّا
يَشْعُرُونَ. وَإِذَا قِيلَ لَهُمْ آمِنُوا كَمَا آمَنَ النَّاسُ قَالُوا
أَنُؤْمِنُ كَمَا آمَنَ السُّفَهَاءُ أَلَا إِنَّهُمْ هُمُ السُّفَهَاءُ
وَلَٰكِن لَّا يَعْلَمُونَ}{البقرہ، \textenglish{11--13}}
\tarjuma{اور جب ان سے کہا جاتا ہے کہ زمین میں فساد نہ پھیلاؤ تو کہتے ہیں: ہم
تو صرف اصلاح کرنے والے ہیں۔ خبردار! یہی لوگ فساد کرنے والے ہیں مگر شعور نہیں
رکھتے۔ اور جب ان سے کہا جاتا ہے کہ ایمان لاؤ جیسے اور لوگ ایمان لائے تو کہتے
ہیں: کیا ہم بھی ایمان لائیں جیسے بے وقوف ایمان لائے؟ خبردار! یہی بے وقوف ہیں
مگر جانتے نہیں۔}

\medskip
\textbf{مزید صفات:} دو رخا پن --- مومنوں سے مل کر ایمان کا دعویٰ، شیاطین
(اپنے سرغنوں) سے مل کر مذاق (\textenglish{14})؛ اللہ کو دھوکا دینے کی کوشش،
حالانکہ دھوکا خود انہی پر پڑتا ہے (\textenglish{9})؛ اور تین علاماتِ نفاق جو
حدیث میں بھی آئیں: جھوٹ بولنا، وعدہ خلافی، اور امانت میں خیانت (بخاری)۔

\medskip
\textbf{کافر و منافق میں فرق:}
\begin{center}
\begin{tabular}{@{}p{3.0cm} p{5.5cm} p{5.5cm}@{}}
\toprule
 & \textbf{کافر} & \textbf{منافق} \\
\midrule
انکار & کھلا اور واضح & زبان سے اقرار، دل میں انکار \\
خطرہ & باہر سے، پہچانا جا سکتا ہے & اندر سے، چھپا ہوا \\
انجام & جہنم & جہنم کا \textbf{سب سے نچلا درجہ} (النساء، \textenglish{145}) \\
\bottomrule
\end{tabular}
\end{center}
\end{hal}

%---------------------------------------------------------------------
\begin{sawal}{سوال \textenglish{3} \ \textnormal{\small(بہار \textenglish{2022}، سوال \textenglish{3} اور سوال \textenglish{4})}}
اہلِ کتاب کے رسول کریم \saw کو پہچاننے اور کتمانِ حق کی مذمت میں آیات کا ترجمہ
و تفسیر لکھیے۔
\end{sawal}

\begin{hal}
\ayat{الَّذِينَ آتَيْنَاهُمُ الْكِتَابَ يَعْرِفُونَهُ كَمَا يَعْرِفُونَ
أَبْنَاءَهُمْ وَإِنَّ فَرِيقًا مِّنْهُمْ لَيَكْتُمُونَ الْحَقَّ وَهُمْ
يَعْلَمُونَ}{البقرہ، \textenglish{146}}
\tarjuma{جن لوگوں کو ہم نے کتاب دی وہ اس (رسول) کو ایسے پہچانتے ہیں جیسے اپنے
بیٹوں کو پہچانتے ہیں، اور ان میں سے ایک گروہ جانتے بوجھتے حق چھپاتا ہے۔}

\smallskip
\textbf{تشریح:} تورات و انجیل میں نبی آخر الزمان \saw کی صفات اتنی واضح تھیں
کہ یہود و نصاریٰ کو اپنی اولاد پہچاننے جیسا یقین تھا --- پھر بھی حسد اور
دنیاوی مفاد کی خاطر حق چھپایا گیا۔ اگلی آیت (\textenglish{147}) فیصلہ سنا دیتی
ہے: \ar{الْحَقُّ مِن رَّبِّكَ فَلَا تَكُونَنَّ مِنَ الْمُمْتَرِينَ} --- حق
تیرے رب کی طرف سے ہے، شک کرنے والوں میں نہ ہو۔

\medskip
\ayat{إِنَّ الَّذِينَ يَكْتُمُونَ مَا أَنزَلَ اللَّهُ مِنَ الْكِتَابِ
وَيَشْتَرُونَ بِهِ ثَمَنًا قَلِيلًا أُولَٰئِكَ مَا يَأْكُلُونَ فِي بُطُونِهِمْ
إِلَّا النَّارَ}{البقرہ، \textenglish{174}}
\tarjuma{جو لوگ اللہ کی نازل کردہ کتاب چھپاتے ہیں اور اس کے بدلے تھوڑی سی
قیمت خریدتے ہیں، وہ اپنے پیٹوں میں سوائے آگ کے کچھ نہیں بھرتے۔}

\smallskip
\textbf{تشریح:} کتمانِ حق کا مالی فائدہ عارضی ہے، مگر انجام آگ ہے --- تشبیہ
سے سمجھایا گیا کہ جو رقم حق چھپا کر کمائی جائے، وہ روزی نہیں بلکہ عذاب ہے۔
آیت \textenglish{175} میں فرمایا: \ar{أُولَٰئِكَ الَّذِينَ اشْتَرَوُا
الضَّلَالَةَ بِالْهُدَىٰ} --- انہوں نے ہدایت کے بدلے گمراہی خرید لی۔
\end{hal}

%---------------------------------------------------------------------
\begin{sawal}{سوال \textenglish{4} \ \textnormal{\small(خزاں \textenglish{2021}، سوال \textenglish{2} اور سوال \textenglish{3})}}
اسلام کے مفہوم اور ''لَا اِکراہ فی الدین'' کی روشنی میں دین میں جبر نہ ہونے
پر آیات کا ترجمہ و تفسیر لکھیے۔
\end{sawal}

\begin{hal}
\ayat{بَلَىٰ مَنْ أَسْلَمَ وَجْهَهُ لِلَّهِ وَهُوَ مُحْسِنٌ فَلَهُ أَجْرُهُ
عِندَ رَبِّهِ وَلَا خَوْفٌ عَلَيْهِمْ وَلَا هُمْ يَحْزَنُونَ}{البقرہ،
\textenglish{112}}
\tarjuma{کیوں نہیں! جو اپنا رخ اللہ کی طرف جھکا دے اور نیکو کار بھی ہو، اس کا
اجر اس کے رب کے پاس ہے، نہ ان پر کوئی خوف ہو گا نہ وہ غمگین ہوں گے۔}

\smallskip
\textbf{تشریح:} یہود و نصاریٰ کے اس دعوے کے جواب میں کہ جنت صرف انہی کے لیے
ہے، اسلام کی حقیقی تعریف دی گئی: \emph{اسلام} نہ کسی نسل کا نام ہے نہ محض
دعویٰ، بلکہ \textbf{اپنا رخ اللہ کی طرف کر دینا} اور ساتھ \textbf{احسان}
(نیکی) بھی۔ صلہ بھی دو ہیں: نہ خوف، نہ غم۔

\medskip
\ayat{لَا إِكْرَاهَ فِي الدِّينِ قَد تَّبَيَّنَ الرُّشْدُ مِنَ الْغَيِّ فَمَن
يَكْفُرْ بِالطَّاغُوتِ وَيُؤْمِن بِاللَّهِ فَقَدِ اسْتَمْسَكَ بِالْعُرْوَةِ
الْوُثْقَىٰ لَا انفِصَامَ لَهَا}{البقرہ، \textenglish{256}}
\tarjuma{دین میں کوئی زبردستی نہیں۔ ہدایت گمراہی سے واضح ہو چکی ہے۔ جو طاغوت
کا انکار کر کے اللہ پر ایمان لائے، اس نے مضبوط کڑے کو تھام لیا جو کبھی ٹوٹنے
والا نہیں۔}

\smallskip
\textbf{تشریح:} دین دل کی رضا سے قبول کی جانے والی چیز ہے، تلوار کے زور سے
نہیں --- کیونکہ حق و باطل کا فرق اتنا واضح ہو چکا ہے کہ زبردستی کی ضرورت ہی
نہیں۔ \emph{طاغوت} یعنی اللہ کے سوا ہر باطل معبود؛ ایمان کو ''مضبوط کڑے''
سے تشبیہ دی گئی جو کبھی نہیں ٹوٹتا۔ یہ آیت اسلام میں \textbf{مذہبی آزادی} کی
بنیادی دلیل ہے۔
\end{hal}

%---------------------------------------------------------------------
\begin{sawal}{سوال \textenglish{5} \ \textnormal{\small(خزاں \textenglish{2021}، سوال \textenglish{1})}}
نماز، زکوٰۃ اور نیکی کی تلقین سے متعلق آیات کا ترجمہ و تفسیر لکھیے، اور بیان
کریں کہ ان میں کن کی اہمیت بیان کی گئی ہے۔
\end{sawal}

\begin{hal}
\ayat{وَأَقِيمُوا الصَّلَاةَ وَآتُوا الزَّكَاةَ وَارْكَعُوا مَعَ
الرَّاكِعِينَ. أَتَأْمُرُونَ النَّاسَ بِالْبِرِّ وَتَنسَوْنَ أَنفُسَكُمْ
وَأَنتُمْ تَتْلُونَ الْكِتَابَ أَفَلَا تَعْقِلُونَ}{البقرہ،
\textenglish{43--44}}
\tarjuma{اور نماز قائم کرو، زکوٰۃ دو اور رکوع کرنے والوں کے ساتھ رکوع کرو۔
کیا تم لوگوں کو نیکی کا حکم دیتے ہو اور اپنے آپ کو بھول جاتے ہو، حالانکہ تم
کتاب پڑھتے ہو؟ کیا تم عقل نہیں رکھتے؟}

\medskip
\textbf{تشریح:} پہلی آیت میں تین احکام: \textbf{نماز} (تعلقِ باللہ)،
\textbf{زکوٰۃ} (تعلقِ مع الخلق)، اور \textbf{جماعت} (''رکوع کرنے والوں کے
ساتھ'')۔ دوسری آیت بنی اسرائیل کے علماء پر تنقید ہے جو دوسروں کو نیکی کی
تلقین کرتے مگر خود عمل نہ کرتے --- یہ اصول \textbf{ہر مبلغ اور معلم} پر لاگو
ہوتا ہے: \emph{قول اور فعل کا تضاد} سب سے بڑی رکاوٹ ہے۔ ''کیا تم عقل نہیں
رکھتے'' ایک تنبیہی سوال ہے، محض استفسار نہیں۔
\end{hal}

%---------------------------------------------------------------------
\begin{sawal}{سوال \textenglish{6} \ \textnormal{\small(بہار \textenglish{2013}، سوال \textenglish{2})}}
حج کے مہینوں اور احرام کی حالت میں ممنوعات پر آیت کا ترجمہ و تفسیر لکھیے۔
\end{sawal}

\begin{hal}
\ayat{الْحَجُّ أَشْهُرٌ مَّعْلُومَاتٌ فَمَن فَرَضَ فِيهِنَّ الْحَجَّ فَلَا
رَفَثَ وَلَا فُسُوقَ وَلَا جِدَالَ فِي الْحَجِّ}{البقرہ، \textenglish{197}}
\tarjuma{حج کے کچھ مقررہ مہینے ہیں، تو جو ان میں حج کی نیت باندھ لے، اس کے
لیے نہ شہوانی گفتگو ہے، نہ نافرمانی، نہ حج میں جھگڑا۔}

\medskip
\textbf{تشریح:} \textbf{اَشْہُرٌ مَعْلُومٰتٌ} سے مراد \emph{شوال، ذی قعدہ،
اور ذی الحجہ کے پہلے دس دن} ہیں --- انہی میں حج کی نیت باندھی جاتی ہے۔ احرام
باندھتے ہی تین چیزیں ممنوع ہو جاتی ہیں: \textbf{رَفَث} (میاں بیوی کا تعلق یا
اس کی گفتگو)، \textbf{فُسُوق} (نافرمانی کی کوئی بھی بات)، اور \textbf{جِدَال}
(بلا وجہ جھگڑا و بحث)۔ آگے فرمایا: \ar{وَتَزَوَّدُوا فَإِنَّ خَيْرَ الزَّادِ
التَّقْوَىٰ} --- زادِ راہ لے لو، اور بہترین زادِ راہ تقویٰ ہے۔ یعنی حج محض
سفری تیاری کا نام نہیں، روحانی پاکیزگی کا نام ہے۔
\end{hal}

%---------------------------------------------------------------------
\begin{sawal}{سوال \textenglish{7} \ \textnormal{\small(بہار \textenglish{2013}، سوال \textenglish{5} اور خزاں \textenglish{2013}، سوال \textenglish{2})}}
طلاق اور خلع کے احکام پر آیات کی روشنی میں نوٹ لکھیے۔
\end{sawal}

\begin{hal}
\ayat{الطَّلَاقُ مَرَّتَانِ فَإِمْسَاكٌ بِمَعْرُوفٍ أَوْ تَسْرِيحٌ
بِإِحْسَانٍ}{البقرہ، \textenglish{229}}
\tarjuma{طلاق دو بار (تک رجوع کی گنجائش کے ساتھ) ہے، پھر یا تو نیکی سے روک
لینا ہے یا بھلائی سے چھوڑ دینا۔}

\medskip
\textbf{تشریح:} پہلی دو طلاقوں کے بعد \textbf{رجوع} (دوبارہ نکاح میں واپسی)
کا حق ہے؛ تیسری طلاق کے بعد نہیں، الا یہ کہ عورت کسی اور جگہ نکاح کرے۔ اسی
آیت میں \textbf{خُلع} کا ذکر ہے: \ar{فَإِنْ خِفْتُمْ أَلَّا يُقِيمَا حُدُودَ
اللَّهِ فَلَا جُنَاحَ عَلَيْهِمَا فِيمَا افْتَدَتْ بِهِ} --- اگر میاں بیوی کو
خدشہ ہو کہ حدودِ الٰہی قائم نہ رکھ سکیں گے تو عورت کے مال دے کر خود کو آزاد
کرا لینے میں کوئی گناہ نہیں۔ فرق: \textbf{طلاق} شوہر کا حق ہے، \textbf{خلع}
بیوی کی درخواست پر، مال کے عوض علیحدگی۔ عدت \textbf{تین حیض} مقرر کی گئی
(\textenglish{228})، تاکہ نسب میں التباس نہ ہو۔
\end{hal}

%---------------------------------------------------------------------
\begin{sawal}{سوال \textenglish{8} \ \textnormal{\small(بہار \textenglish{2022}، سوال \textenglish{7} اور خزاں \textenglish{2013}، سوال \textenglish{4})}}
قسموں کے احکام کی روشنی میں آیات کا ترجمہ و تفسیر لکھیے۔
\end{sawal}

\begin{hal}
\ayat{وَلَا تَجْعَلُوا اللَّهَ عُرْضَةً لِّأَيْمَانِكُمْ أَن تَبَرُّوا
وَتَتَّقُوا وَتُصْلِحُوا بَيْنَ النَّاسِ}{البقرہ، \textenglish{224}}
\tarjuma{اور اللہ کے نام کو اپنی قسموں کا ایسا نشانہ نہ بناؤ کہ (اس کے بہانے)
نیکی، تقویٰ اور لوگوں کے درمیان صلح سے رک جاؤ۔}

\smallskip
\ayat{لَا يُؤَاخِذُكُمُ اللَّهُ بِاللَّغْوِ فِي أَيْمَانِكُمْ وَلَٰكِن
يُؤَاخِذُكُم بِمَا كَسَبَتْ قُلُوبُكُمْ}{البقرہ، \textenglish{225}}
\tarjuma{اللہ تمہیں تمہاری لغو قسموں پر نہیں پکڑتا، بلکہ اس پر پکڑتا ہے جو
تمہارے دلوں نے کمایا ہو۔}

\medskip
\textbf{تشریح:} پہلی آیت میں منع کیا گیا کہ اللہ کے نام کی قسم کھا کر نیکی کا
کام چھوڑ دیا جائے (مثلاً ''میں نے قسم کھائی، اب صلح نہیں کروں گا'') --- ایسی
قسم توڑ کر نیکی کرنا ہی افضل ہے۔ دوسری آیت \textbf{لغو قسم} (بلا ارادہ، عادتاً
منہ سے نکل جانے والی، جیسے ''نہیں، خدا کی قسم'') اور \textbf{قصدی قسم} میں
فرق کرتی ہے --- صرف قصدی، دل کے پختہ ارادے سے کھائی گئی قسم پر مؤاخذہ ہے۔
قصدی قسم توڑنے پر کفارہ ہے: دس مسکینوں کو کھانا یا کپڑا، یا ایک غلام آزاد
کرنا؛ استطاعت نہ ہو تو تین روزے (المائدہ، \textenglish{89})۔
\end{hal}

%---------------------------------------------------------------------
\begin{sawal}{سوال \textenglish{9} \ \textnormal{\small(بہار \textenglish{2013}، سوال \textenglish{4})}}
سورۃ البقرہ کی روشنی میں صبر کی تعریف اور فضیلت، نیز شہید کے فضائل و احکام
بیان کریں۔
\end{sawal}

\begin{hal}
\ayat{يَا أَيُّهَا الَّذِينَ آمَنُوا اسْتَعِينُوا بِالصَّبْرِ وَالصَّلَاةِ
إِنَّ اللَّهَ مَعَ الصَّابِرِينَ}{البقرہ، \textenglish{153}}
\tarjuma{اے ایمان والو! صبر اور نماز سے مدد لو، بے شک اللہ صبر کرنے والوں کے
ساتھ ہے۔}

\smallskip
\textbf{تعریف۔} صبر یعنی مصیبت پر بے قراری کے بجائے ثابت قدمی، زبان سے
شکایت نہ کرنا اور اللہ کی رضا پر راضی رہنا۔ تین قسمیں: اطاعت پر صبر، گناہ سے
رکنے پر صبر، اور مصیبت پر صبر۔

\medskip
\textbf{فضیلت۔} \ar{وَبَشِّرِ الصَّابِرِينَ الَّذِينَ إِذَا أَصَابَتْهُم
مُّصِيبَةٌ قَالُوا إِنَّا لِلَّهِ وَإِنَّا إِلَيْهِ رَاجِعُونَ} ---
صبر کرنے والوں کو خوشخبری دیجیے جو مصیبت پر ''انا للہ وانا الیہ راجعون''
کہتے ہیں (\textenglish{155--156})۔ آگے فرمایا: \ar{أُولَٰئِكَ عَلَيْهِمْ
صَلَوَاتٌ مِّن رَّبِّهِمْ وَرَحْمَةٌ} --- انہی پر ان کے رب کی رحمتیں اور
برکتیں ہیں (\textenglish{157})۔

\medskip
\textbf{شہید کے فضائل۔} \ar{وَلَا تَقُولُوا لِمَن يُقْتَلُ فِي سَبِيلِ
اللَّهِ أَمْوَاتٌ بَلْ أَحْيَاءٌ وَلَٰكِن لَّا تَشْعُرُونَ} --- اللہ کی راہ
میں مارے جانے والوں کو مردہ مت کہو، بلکہ وہ زندہ ہیں مگر تمہیں شعور نہیں
(\textenglish{154})۔ شہید کے احکام: غسل نہیں دیا جاتا، اسی خون آلود لباس میں
دفن کیا جاتا ہے، اور جنازہ حنفیہ کے نزدیک پڑھا جاتا ہے۔
\end{hal}
